% Requires graphicx and the manuscript's existing natbib configuration.
% Citation keys are already present in related_work/conv_related_work.bib.
\section{Introduction}
\label{sec:introduction}

Long-context language models support tasks that require reasoning over
extended documents, retrieving distant evidence, and processing large code
contexts. Their practical use depends on processing long prompts efficiently
before autoregressive generation begins. During this prefilling stage,
attention evaluates a quadratic number of query--key interactions, making
long prompts expensive even when model parameters remain unchanged.
Memory-efficient kernels such as FlashAttention~\citep{dao2022flashattention}
reduce the cost of executing exact attention, while sparse attention reduces
computation by omitting selected interactions. For pretrained language models,
the central challenge is to identify which interactions can be skipped without
discarding information needed for prediction.

Block-sparse prefilling addresses this challenge by selecting groups of
query--key interactions that can be processed efficiently by sparse kernels.
Recent methods construct input-dependent patterns using structured attention
families~\citep{jiang2024minference}, adapt patterns and budgets to individual
inputs and heads~\citep{lai2025flexprefill}, or estimate block importance from
antidiagonal samples~\citep{xu2025xattention}. Learned selectors such as
SeerAttention~\citep{gao2025seerattention} instead train lightweight gates to
predict block importance. These advances make sparse prefilling practical,
but leave a question about selection quality: how well does an estimated
block score reflect the benefit of retaining that block within the selected
set? An inexpensive importance proxy need not order blocks in the same way
as their contributions to the final prediction.

We investigate this distinction through two observations. First, in a
controlled replacement experiment, a block excluded by the initial ranking
can reduce the ground-truth answer loss when substituted for a higher-scored
selected block, with the block budget and all other selections held fixed.
Thus, a higher initial score does not always imply greater prediction
utility. Second, selected-block masks exhibit vertical and diagonal
structures, suggesting that nearby scores may provide useful context.
Matching candidates with similar center scores and examining residual
associations show that neighborhood scores contain additional information
about utility in some examined layers. This relationship varies across
layers and is observational rather than causal. Together, these findings
motivate learning how local score structure should influence block selection.

We propose \emph{Conv}, a learnable convolutional refinement of block scores
for sparse prefilling. Figure~\ref{fig:overall_framework} places the method
within the language-model inference pipeline. Starting from RoPE-encoded
queries and keys, an antidiagonal estimator produces an initial block-score
map. A layer- and head-specific $7\times7$ convolution then combines each
score with its local neighborhood before block selection. The selected
indices guide attention computed from the original-resolution queries,
keys, and values. The convolution therefore modifies the selection criterion;
it does not replace the attention weights computed within retained blocks.
Only the convolutional kernels are optimized, using attention statistics
from the frozen model and task-relevant supervision. Their parameters can
adapt the contribution of neighboring scores to each layer and head, while
the pretrained language-model parameters remain frozen and standard decoding
is unchanged.

\begin{figure*}[t]
    \centering
    \includegraphics[width=\textwidth]{Figs/overall.png}
    \caption{
        Position of Conv within long-context language-model inference.
        During prefilling, antidiagonal scoring produces initial block scores
        from RoPE-encoded queries and keys. A learnable $7\times7$ convolution
        refines these scores before block selection; top-$k$ selection is
        illustrated. The selected indices guide block-sparse attention using
        full-resolution queries, keys, and values. Only the convolutional
        kernels are trained. The language-model parameters remain frozen,
        and subsequent autoregressive decoding follows the standard path.
    }
    \label{fig:overall_framework}
\end{figure*}

We evaluate Conv on Qwen3-8B and Llama-3.1-8B using RULER and LongBench.
Fixed-budget top-$k$ comparisons isolate the effect of score refinement,
whereas top-$p$ efficiency measurements examine adaptive selection when the
retained block count depends on the input. On Llama-3.1-8B, Conv improves
RULER accuracy over the Xattn baseline across all tested combinations of
context length and keep ratio. The broader results also reveal dependence
on the model and operating point: the more aggressive Qwen3 configuration
does not uniformly improve benchmark scores, and short prompts do not
benefit from the measured sparse-attention speedups. These comparisons
establish the scope of the gains without assuming that a common threshold
implies equal sparsity or equal accuracy.

Our contributions are threefold:
\begin{itemize}
    \item We distinguish estimated block importance from prediction utility
    through controlled block replacements and identify layer-dependent
    neighborhood information that supplements a block's own score.

    \item We introduce a lightweight convolutional score-refinement module
    that learns local selection cues for each layer and head while keeping
    the language model frozen.

    \item We evaluate selection quality and efficiency under complementary
    top-$k$ and top-$p$ protocols. On Llama-3.1-8B, Conv gains an average of
    3.07 RULER points over Xattn across 12 matched-budget settings. In a
    separate top-$p=0.9$ benchmark at 128K, its prefill attention path reaches
    $3.67\times$ speedup over full attention, compared with $2.87\times$
    for Xattn.
\end{itemize}
